\chapter{Instruction Syntax}
\label{ch:instructions}

\section{Anatomy of an Instruction}

An instruction is written as a mnemonic, an optional status update source (\texttt{<SUS>}), an optional condition
(\texttt{<COND>}), and then its operands:

\begin{center}
	\texttt{<MNEMONIC>\{[<SUS>]\}\{|<COND>\} <OPERAND>, ...}
\end{center}

The three parts must be written in that order and with no spaces between them. The operands follow after at least one
space.

% qa:assemble
\begin{lstlisting}
    ADDI r1, #1                  ; mnemonic and operands
    ADDI[N] r1, #1               ; with a status update source
    ADDI|== r1, #1               ; with a condition
    ADDI[N]|== r1, #1            ; with both
\end{lstlisting}

\section{Status Update Source}
\label{sec:status-update-syntax}

\texttt{[N]}, \texttt{[A]} or \texttt{[S]} after the mnemonic selects what updates the status flags: nothing, the ALU
result, or the shifter result. Without it the instruction uses its default, which is \texttt{[A]} for most ALU
instructions. The default of each instruction and the meaning of each source are in the \archmanual{}.

Only ALU instructions have a status update source. \mnemonic{MOVI} and \mnemonic{MOVR} accept \texttt{[A]} and
\texttt{[N]} but not \texttt{[S]}, because they have no shifted operand. Memory and control-flow instructions never
update the flags and must not be written with a status update source.

\section{Conditions}

\texttt{|<COND>} after the mnemonic, and after the status update source if there is one, makes the instruction
conditional. The sixteen conditions are:

\begin{table}[H]
	\centering
	\begin{tabular}{lp{9cm}}
		\toprule
		\textbf{Written as}                                    & \textbf{Meaning}                                  \\
		\midrule
		\texttt{AL}                                            & Always (the default when no condition is written) \\
		\texttt{==}, \texttt{!=}                               & Equal, not equal                                  \\
		\texttt{CS}, \texttt{CC}                               & Carry set, carry clear                            \\
		\texttt{NS}, \texttt{NC}                               & Negative set, negative clear                      \\
		\texttt{OS}, \texttt{OC}                               & Overflow set, overflow clear                      \\
		\texttt{HI}, \texttt{LO}                               & Unsigned higher, unsigned lower or same           \\
		\texttt{>=}, \texttt{<{}<}, \texttt{>{}>}, \texttt{<=} & Signed comparisons                                \\
		\texttt{NV}                                            & Never                                             \\
		\bottomrule
	\end{tabular}
	\caption{Condition Codes in Source}
	\label{tab:conditions-in-source}
\end{table}

The Condition Codes chapter of the \archmanual{} defines the flags each condition tests and how to choose between them.

\section{Choosing the Immediate Format}

The assembler picks the encoding from how the operands are written. For the ALU instructions that have both forms, an
immediate operand on its own is encoded in the 16-bit immediate format, and an immediate operand followed by a shift is
encoded in the 8-bit short immediate format:

% qa:assemble
\begin{lstlisting}
    ADDI r1, #5                  ; immediate format, 16-bit value
    ADDI r1, #5, LSL #1          ; short immediate format, 8-bit value
\end{lstlisting}

The register format is selected by the mnemonic: write \mnemonic{ADDR} with register operands, not \mnemonic{ADDI}.

\section{Mnemonics}

The assembler recognizes the mnemonics in Table~\ref{tab:mnemonic-families}. Each is documented in the chapter of the
\archmanual{} named in the table.

\begin{table}[H]
	\centering
	\small
	\begin{tabularx}{\textwidth}{lXp{3.3cm}}
		\toprule
		\textbf{Family}      & \textbf{Mnemonics}                                                                                                                                                                                                                                                     & \textbf{\archmanual{} Chapter}      \\
		\midrule
		ALU                  & \mnemonic{ADDI}/\mnemonic{ADDR}, \mnemonic{ADCI}/\mnemonic{ADCR}, \mnemonic{SUBI}/\mnemonic{SUBR}, \mnemonic{SBCI}/\mnemonic{SBCR}, \mnemonic{ANDI}/\mnemonic{ANDR}, \mnemonic{ORRI}/\mnemonic{ORRR}, \mnemonic{XORI}/\mnemonic{XORR}, \mnemonic{MOVI}/\mnemonic{MOVR} & ALU Instructions                    \\
		ALU test             & \mnemonic{CMNI}/\mnemonic{CMNR}, \mnemonic{CNCI}/\mnemonic{CNCR}, \mnemonic{CMPI}/\mnemonic{CMPR}, \mnemonic{CPCI}/\mnemonic{CPCR}, \mnemonic{TSAI}/\mnemonic{TSAR}, \mnemonic{TSOI}/\mnemonic{TSOR}, \mnemonic{TSXI}/\mnemonic{TSXR}                                  & ALU Instructions                    \\
		Memory               & \mnemonic{LOAD}, \mnemonic{STOR}                                                                                                                                                                                                                                       & Memory Access Instructions          \\
		Addresses            & \mnemonic{LDEA}, \mnemonic{LDEL}, \mnemonic{LJMP}, \mnemonic{LJSR}                                                                                                                                                                                                     & Control Flow Instructions           \\
		Branches and returns & \mnemonic{BRAN}, \mnemonic{BRSR}, \mnemonic{RETS}                                                                                                                                                                                                                      & Control Flow Instructions           \\
		Coprocessor          & \mnemonic{COPI}, \mnemonic{COPR}, \mnemonic{EXCP}, \mnemonic{WAIT}, \mnemonic{RETE}, \mnemonic{RSET}, \mnemonic{ETFR}, \mnemonic{ETTR}                                                                                                                                 & Coprocessor Instructions            \\
		Maths and DMA        & \mnemonic{MULU}, \mnemonic{MULS}, \mnemonic{DIVU}, \mnemonic{DIVS}, \mnemonic{DMAR}, \mnemonic{DMAW}, \mnemonic{DMAT}                                                                                                                                                  & Coprocessor Instructions            \\
		Other meta           & \mnemonic{NOOP}, \mnemonic{SHFT}                                                                                                                                                                                                                                       & Meta-Instructions, ALU Instructions \\
		\bottomrule
	\end{tabularx}
	\caption{Mnemonics Recognized by the Assembler}
	\label{tab:mnemonic-families}
\end{table}

Many of these are \textit{meta-instructions}: names for a single real instruction encoding. The Meta-Instructions
chapter of the \archmanual{} lists what each one assembles to. The assembler has no pseudoinstructions that expand to
more than one instruction.
